\documentclass[11pt,a4paper]{article}

\usepackage[T1]{fontenc}
\usepackage[margin=2.3cm]{geometry}
\usepackage{booktabs}
\usepackage{amsmath}
\usepackage{array,tabularx}
\usepackage{enumitem}
\usepackage{titlesec}
\usepackage{parskip}
\usepackage{caption}
\usepackage{xcolor}
\usepackage{hyperref}

\hypersetup{colorlinks=true,linkcolor=blue,citecolor=blue,urlcolor=blue}
\captionsetup[table]{position=top,skip=4pt,labelfont=bf}
\newcolumntype{Y}{>{\raggedright\arraybackslash}X}
\setlist{itemsep=3pt,topsep=4pt}
\titlespacing*{\section}{0pt}{11pt}{5pt}
\titlespacing*{\subsection}{0pt}{8pt}{3pt}
\newcommand{\field}[1]{\textbf{#1:}~}
\newcommand{\hz}[1]{{\color{blue}[#1 --- HZ]}}
\setlength{\emergencystretch}{2em}

\title{\textbf{NR Sidelink Measurement Plan}\\[0.4em]
\large Reliable Communication Between Agricultural Machines}
\author{}
\date{Study plan: 26 October 2026}

\begin{document}
\maketitle

\section{Study goal}

The goal is to find out how well NR sidelink Mode 2 can support communication
between agricultural machines. We want to know when messages arrive on time,
when communication becomes unreliable, and what causes the change.

The study will answer five questions:
\begin{enumerate}
  \item \textbf{Can we trust the measurements?}
        Do the communication procedures, packet counts, and timing checks
        give the expected results?
  \item \textbf{How well are messages delivered?}
        How many arrive before their deadline, how long do they take,
        and how long can a machine go without a useful update?
  \item \textbf{Where does communication work well?}
        How do distance, movement, blocked paths, and the number of machines
        change delivery?
  \item \textbf{What helps, and what does it cost?}
        How much do repeated transmissions and delivery feedback help?
        How much extra delay and radio resource do they require?
  \item \textbf{What does this mean for farm work?}
        Under which conditions can the communication needs of a harvester,
        grain cart, or group of machines be met?
\end{enumerate}

The main outcome will be a clear account of the conditions under which each
application meets its communication needs. Useful data rate will also be
measured, especially when many machines are active.

This is a simulation study. Physical radio measurements can later test how
well its findings carry over to actual equipment.

\section{Measurements}
\label{sec:metrics}

Before each experiment, state the message size, how often messages are
created, who should receive them, and the delivery deadline. Set these from
the application's needs before judging the result.

\begin{table}[htbp]
\centering
\caption{Measurements and their meaning}
\label{tab:metrics}
\small
\begin{tabularx}{\textwidth}{lY}
\toprule
\textbf{Measurement} & \textbf{What it tells us} \\
\midrule
Delivery before the deadline &
The share of expected message deliveries that happen on time. This is the
main measure of reliability. \\
Message delay &
Time from message creation at the sending application to delivery at the
receiving application. \\
Delay variation &
How much delay changes between consecutive messages. \\
Gap between updates &
How long the receiver waits between successful updates. For coordination,
also measure the gap between fresh updates that arrive on time. \\
Loss bursts &
How many messages in a row miss their deadline. \\
Useful data rate &
Application data delivered per second, counting each message once at each
intended receiver. \\
Radio resource use &
How much of the shared radio resource is used, and how much is needed to
deliver one useful message. \\
\bottomrule
\end{tabularx}
\end{table}

Two simple calculations are central:
\begin{equation}
\text{On-time delivery (\%)} =
100 \times
\frac{\text{Expected deliveries completed before the deadline}}
     {\text{All expected deliveries}}.
\label{eq:delivery}
\end{equation}
\begin{equation}
\text{Message delay} =
\text{Application delivery time} - \text{Application creation time}.
\label{eq:delay}
\end{equation}

For example, if 1,000 messages should reach one machine and 990 arrive before
their deadline, on-time delivery is 99\%. If a message should reach several
machines, count its expected delivery to each machine separately.

Report typical delay and the slower end of the delay distribution. The median
is the middle delay. The 95th and 99th percentiles are the delays within which
95\% and 99\% of delivered messages arrive. Show these beside the delivery
percentage, because delay values alone leave out messages that never arrive.

Also report what share of all expected deliveries is complete by each elapsed
time. If 99\% have not arrived by the end of observation, say that the
all-message 99th percentile was not reached.

For delay variation, compare consecutive generated messages only when both
arrive. Keep missing messages visible in the loss and gap results.

\section{Experiments}
\label{sec:experiments}

Each experiment has one main goal. Begin with simple cases, then add distance,
movement, repeated transmissions, and more machines.

\begin{table}[htbp]
\centering
\caption{Experiment overview}
\small
\begin{tabularx}{\textwidth}{lYY}
\toprule
\textbf{Test} & \textbf{Goal} & \textbf{Main result} \\
\midrule
M0: Measurement checks & Check that results can be trusted &
Evidence that timing, delivery, and radio procedures behave as expected. \\
M1: Simple link & Establish a reference &
Delay, delivery, and useful data rate under easy conditions. \\
M2: Distance & Find the working range &
Delivery and delay as machines move farther apart. \\
M3: Movement & Find the effect of speed and changing position &
Delivery, delay, and update gaps during movement. \\
M4: Repeated transmissions & Find the benefit and cost of retries &
Delivery improvement, added delay, and extra resource use. \\
M5: More machines & Find how much traffic can be supported &
The point at which the application target stops being met. \\
M6: Farm tasks & Check application needs &
Working conditions for selected agricultural tasks. \\
\bottomrule
\end{tabularx}
\end{table}

\subsection{M0: Check that the measurements can be trusted}

\field{Goal} Confirm that the basic measurements and communication procedures
are correct before drawing performance conclusions.

\field{Test} Start with cases whose answers are already known. Check that
every created message is accounted for, a known delivery time gives the
correct delay, and a deliberately lost message appears as a loss. Check
received power, noise, and available space for data against simple calculations.

Then check the main communication steps: finding a resource, announcing a
transmission, receiving data, handling repeated copies, sending delivery
feedback, and dropping messages after their allowed waiting time.
Include a case where a machine is transmitting when it should receive.

Compare the relevant procedures with independent examples and the applicable
sidelink rules~\cite{ts38214}. Test reception across a range of signal
conditions. Separate failures to read the control information from failures
to read the message data.

Where a simplified calculation predicts detailed reception results, check
its accuracy on separate cases that were not used to adjust it. Include
interference and repeated transmissions. Set the allowed prediction error
before making the comparison.

\field{Result} A record of what passed, what failed, and the conditions covered
by each check. Correct problems affecting later experiments before continuing.

\subsection{M1: Establish the simple-link reference}

\field{Goal} Measure communication between two machines under easy conditions.

\field{Test} Place two stationary machines close together with a clear path
between them. Send messages at the application's normal rate. Measure delay,
on-time delivery, update gaps, and useful data rate.

Run a separate test with data always waiting to be sent. This shows the
rate the tested communication arrangement can sustain. Record queue growth
and radio resource use as well as delivered data.

\field{Result} A reference for later comparisons. Show where time is spent:
preparing a message, waiting to send, completing transmission and any retries,
and delivering it to the receiving application.

Keep the normal-message test and the maximum-load test separate. They answer
different questions.

\subsection{M2: Find how distance changes communication}

\field{Goal} Find the distances over which the application target can be met.

\field{Test} Keep the machines stationary and increase their separation.
Compare a clear path with a blocked or obstructed path. Keep the remaining
settings the same.

Use an initial trial to locate the distance where delivery starts to fall.
Then choose the full distance range and closer spacing around that change
before collecting the main results. Describe any obstruction clearly;
a general obstructed-path case does not by itself represent a machine body.

\field{Result} Plot on-time delivery, delay, update gaps, and useful data rate
against distance. For a stated deadline, report the tested distances meeting
90\% and 99\% delivery, or the application's own target.

Show the last passing distance and the next distance where the target is
missed or the evidence is unclear. If the target is still met at the farthest
tested point, report that tested distance as a lower bound on range.

A range based only on receiving one transmission successfully must be reported
separately from the range for delivering application messages on time.

\subsection{M3: Find how movement changes communication}

\field{Goal} Find whether movement causes extra delay, losses, or long gaps
between useful updates.

\field{Test} Begin with one stationary machine and one machine moving past it.
Repeat at speeds relevant to the farm task. Include an approach and a departure.
Then test both machines moving, with their directions clearly stated.

Compare moving and stationary cases at similar distances. Repeat complete
passes and record how long each pass spends in each distance range.

\field{Result} Show delivery, delay, and update gaps against both distance and
speed. Identify where movement changes performance and whether the receiver
can still follow the changing radio signal.

If a distance range has too few observations, show that limitation instead
of giving it the same weight as a well-observed range.

\subsection{M4: Find whether repeated transmissions are worthwhile}

\field{Goal} Measure how much repeated transmissions improve delivery and what
they add in delay and radio resource use.

\field{Test} Compare three cases:
\begin{enumerate}
  \item Send the message once.
  \item Send a planned repeat without waiting for delivery feedback.
  \item Repeat when delivery feedback shows that another attempt is needed.
\end{enumerate}

Make this comparison for the same sender, intended receiver, message traffic,
and radio conditions. Give the two repeating methods the same limit on the
number of attempts. Choose easy, moderate, and difficult distances during
the initial trial, before collecting the main results.

First keep the radio resource layout the same to compare the retry methods.
Then include the change in resources needed for feedback to measure the
overall benefit and cost. Study repeats to several receivers separately.

\field{Result} Show the gain in on-time delivery beside the change in delay,
number of attempts, feedback failures, and resources used per delivered
message. Check whether a retry arrives early enough to help the application.

\subsection{M5: Find how many machines and messages can be supported}

\field{Goal} Find how performance changes as more machines share the channel.

\field{Test} Keep the field area fixed and increase the number of machines.
First hold each machine's message rate constant. Then increase the message
rate and test a case where many machines send an event message at the same
time.

Repeat the test with new machine positions and radio conditions. Compare
machines at similar distances, and record how busy the channel appears to
each machine.

\field{Result} Show on-time delivery, delay, update gaps, useful data rate,
waiting queues, and resource use against machine count and message load.
Report results for individual machines as well as the overall average.

Investigate losses caused by waiting too long, being unable to receive while
transmitting, and failures to read control information or data. Use transmission
records to check whether overlapping transmissions explain the failures.

Report the largest tested machine count and message load that meet the
application target with enough confidence. State the field size, traffic,
and radio conditions beside that result.

\subsection{M6: Check the communication needs of farm tasks}

\field{Goal} Connect the measurements to the needs of agricultural applications.

\field{Test} Study the following tasks:
\begin{table}[htbp]
\centering
\caption{Agricultural tasks and their communication questions}
\small
\begin{tabularx}{\textwidth}{YY}
\toprule
\textbf{Task} & \textbf{Question} \\
\midrule
Harvester and grain cart &
Do coordination messages arrive on time as separation and relative position
change? How long can the gap between useful updates become? \\
Several machines working together &
What happens during parallel passes, turns at the field edge, and messages
sent by several machines at once? \\
Messages to a field gateway &
Can machines send updates reliably to another sidelink device acting as a
gateway? How do feedback and packet recovery change delay and delivery? \\
\bottomrule
\end{tabularx}
\end{table}

For each task, first state the message size, message rate, deadline, delivery
target, and longest acceptable update gap. Include the required distances,
speeds, and number of machines.

Follow complete task segments, including turns or position changes that matter
to communication. Identify whether the paths are recorded machine movements
or constructed examples.

When comparing packet recovery above the radio layer, keep radio retries
unchanged. When comparing radio retries, keep the other recovery settings
unchanged. Measure delivery at the application so that waiting, recovery,
and reassembly are included.

\field{Result} A table showing which combinations of distance, speed, machine
count, and message load meet each task's needs, miss them, or need more evidence.

\section{Fair comparisons}

\begin{itemize}
  \item \textbf{Change one main condition at a time.}
        After understanding its effect, test important combinations such as
        distance with obstruction, speed with transmission rate settings,
        and machine count with message load.
  \item \textbf{Compare the same situation.}
        Keep machine paths, generated messages, and outside radio conditions
        matched when comparing methods. A shared random seed alone does not
        guarantee this.
  \item \textbf{Count every expected delivery.}
        Decide the intended receivers when each message is created. Keep
        failed deliveries in the results, including messages that expire
        before being sent. Count repeated copies only once.
  \item \textbf{Give the final messages enough time.}
        Continue normal traffic while the final measured messages reach
        their deadlines. Set an extra observation period for late arrivals
        before the test begins, and report messages still undelivered at its end.
  \item \textbf{Keep timing definitions consistent.}
        State where every delay measurement starts and ends. Record added
        processing delays explicitly.
  \item \textbf{Keep data-rate counts clear.}
        Separate unique transmitted data from copies delivered to multiple
        receivers. Record the time interval used to calculate each rate.
\end{itemize}

Set the application targets explicitly. Meeting an on-time delivery target
does not by itself establish every standardized quality-of-service
requirement~\cite{ts23287,ts23501}.

\section{Deciding whether a target is met}

Begin with a small trial to choose sensible test distances, run lengths,
settling time, and numbers of repeats. Check that queues and resource choices
have settled before counting the main measurements. Set the final experiment
conditions and analysis rules before the main study. Keep the trial data
separate.

Repeat whole experiments with new positions, radio conditions, or movement
passes. Messages from one run often share the same conditions. Deliveries of
one broadcast message to many receivers also share conditions. Account for
these links when estimating uncertainty; they are not all independent repeats.

Show a 95\% confidence interval with the main results. This gives a range
showing the uncertainty in the estimate. For delivery targets:
\begin{itemize}
  \item \textbf{Met:} the whole interval is at or above the target.
  \item \textbf{Missed:} the whole interval is below the target.
  \item \textbf{Unclear:} the interval includes both passing and failing values.
\end{itemize}

For a 99\% target, an interval of 99.1--99.6\% supports a pass.
An interval of 98.4--98.9\% supports a miss. An interval of 98.8--99.2\%
leaves the decision unclear.

For network results, estimate uncertainty by repeating or resampling whole
runs, keeping the messages within each run together. For independent reception
trials, use intervals appropriate for counts of successes and failures~\cite{nist}.
State the calculation used, including whether an interval is one-sided or
two-sided. Use the same rule throughout each comparison.

A short test with no observed losses does not prove perfect delivery.
Very high reliability targets need enough independent evidence. If the
planned measurements cannot resolve a target, report the result as unclear.
Do not keep adding runs only until the desired answer appears.

When reporting a range or capacity limit, show the neighboring tested points
and their uncertainty. Distinguish confidence at individual points from
confidence in an entire range.

\section{Records and results}

Keep a record of the experiment conditions, machine positions and movements,
message creation and delivery times, intended receivers, transmission attempts,
feedback, and reasons for missed delivery. Record failures as carefully as
successes. Save the study version and enough information to repeat each test.

Check that all generated messages are accounted for. Report messages that
arrived on time, arrived late, were dropped, or were still undelivered when
observation ended. Mark interrupted tests separately.

For each experiment, report:
\begin{enumerate}
  \item The question and the conditions tested.
  \item Delivery percentage, delay, update gaps, and useful data rate.
  \item Differences between repeated runs and between machines.
  \item The confidence interval and target decision.
  \item The observations that explain the result.
\end{enumerate}

Present the main findings in this order: measurement checks, simple-link
reference, distance, movement, repeated transmissions, machine count and
load, and farm tasks. Produce the tables and figures from the saved records
so every reported result can be checked.

\section{Order of work}

\begin{enumerate}
  \item Agree on the application questions and delivery targets.
  \item Complete the basic measurement checks in M0.
  \item Run the small trial and set the final conditions and repeat counts.
  \item Finish the remaining M0 checks under those conditions.
  \item Complete M1 and M2, followed by M3 and M4, then M5 and M6.
  \item Check the records, calculate uncertainty, and judge each target.
  \item Explain where communication meets the application's needs and where
        it needs improvement.
\end{enumerate}

\begin{thebibliography}{9}

\bibitem{ts38214}
3GPP TS 38.214 V16.17.0, \emph{NR; Physical layer procedures for data},
Release 16, clauses 8.1.3--8.1.4.
\url{https://www.etsi.org/deliver/etsi_ts/138200_138299/138214/16.17.00_60/ts_138214v161700p.pdf}.

\bibitem{ts23287}
3GPP TS 23.287 V19.0.0, \emph{Architecture enhancements for 5G System (5GS)
to support Vehicle-to-Everything (V2X) services}, clauses 5.4.3--5.4.4.
\url{https://www.etsi.org/deliver/etsi_ts/123200_123299/123287/19.00.00_60/ts_123287v190000p.pdf}.

\bibitem{ts23501}
3GPP TS 23.501 V18.8.0, \emph{System architecture for the 5G System (5GS)},
clauses 5.7.3.4--5.7.3.5.
\url{https://www.etsi.org/deliver/etsi_ts/123500_123599/123501/18.08.00_60/ts_123501v180800p.pdf}.

\bibitem{nist}
NIST, \emph{Exact Binomial Confidence Limits}, Dataplot Reference Manual.
\url{https://itl.nist.gov/div898/software/dataplot/refman2/auxillar/exacbici.htm}.

\end{thebibliography}
\end{document}
